\section{Shell Sort}
O Shell Sort é um algoritmo de ordenação que expande as ideias do \textit{Inserção} para melhorar sua eficiência em determinados casos. Ele organiza os elementos de uma lista ao permitir a troca de itens que estão distantes uns dos outros, utilizando um conceito chamado de h-ordenação. A ideia principal é dividir a lista em sublistas baseadas em um intervalo $h$ e ordenar esses subconjuntos de forma independente. Com isso, os elementos são gradualmente movidos para posições mais próximas de sua ordenação final. À medida que o intervalo $h$ diminui, o algoritmo refina a ordenação até que todas as sublistas sejam reduzidas a elementos consecutivos, momento em que o método converge para um \textit{Inserção} final mais eficiente.

\lstinputlisting[language=C]{funcoes-c/f_shell_sort.txt}
% \begin{algorithm}
%     \caption{ShellSort}\label{alg-shell-sort:cap}
%     \DontPrintSemicolon
%     \Entrada{um vetor $v$ de $n$ elementos}
%     \Resultado{o vetor $v$ em ordem crescente}
%     \Inicio{
%         $gap \gets n / 2$\;
%         \Enqto{$gap > 0$}{
%             \Para{$i = gap$ \Ate $\,n - 1\,$}{
%                 $chave \gets v[i]$\;
%                 $j \gets i - gap$\;
%                 \Enqto{$j \geq 0$ \textbf{e} $v[j] > chave$}{
%                     $v[j+gap] \gets v[j]$\;
%                     $j \gets j - gap$\;
%                 }
%                 $v[j+gap] \gets chave$\;
%             }
%             $gap \gets gap / 2$\;
%         }
%     }
% \end{algorithm}

\subsection{Prova de correção}
Para o caso particular em que $n = 1$, o vetor está trivialmente ordenado, e nenhuma movimentação é realizada. Já no caso geral, em que $n > 1$, a variável $gap$ sempre atinge o valor $1$. Nesse ponto, os mesmos passos do algoritmo \textit{Inserção} são executados, garantindo a ordenação do vetor. Portanto, o algoritmo é correto.

\subsection{Complexidade de tempo e espaço}
TODO: Encontrar e citar bibliografias que sustentem as seguinte informações: \textbf{melhor caso:} linear, \textbf{caso médio:} $O(n^{1.5}) \approx O(n \cdot log_2(n))$ e \textbf{pior caso:} $O(n^2)$.
